% from ega3-1-fr.tex
\oldpage[0\textsubscript{III}]{54}
qui donne en passant à la cohomologie des diagrammes commutatifs
\[
\label{0.12.1.5.4-fr}
\tag{12.1.5.4}
\xymatrix@C=7pt{
H^p(X,\mathcal F)\otimes_A H^q(X,\mathcal G)\ar[r]\ar[d]
&H^{p+q}(X,\mathcal H)\ar[d]\\
H^p(\Gamma(X',\psi^*(\mathcal L^\bullet)))
\otimes_{\Gamma(X',\psi^*(\mathcal O_X))}
H^q(\Gamma(X',\psi^*(\mathcal M^\bullet)))\ar[r]
&H^{p+q}(\Gamma(X',\psi^*(\mathcal N^\bullet)))
}
\]
Mais comme $\psi^*(\mathcal L^\bullet)$, $\psi^*(\mathcal M^\bullet)$ et
$\psi^*(\mathcal N^\bullet)$ sont des \emph{résolutions} de
$\psi^*(\mathcal F)$, $\psi^*(\mathcal G)$, $\psi^*(\mathcal H)$
respectivement, on a un diagramme commutatif (G, II, 6.6.1)
\[
\label{0.12.1.5.5-fr}
\tag{12.1.5.5}
\xymatrix@C=7pt{
H^p(\Gamma(X',\psi^*(\mathcal L^\bullet)))
\otimes_{\Gamma(X',\psi^*(\mathcal O_X))}
H^q(\Gamma(X',\psi^*(\mathcal M^\bullet)))\ar[r]\ar[d]
&H^{p+q}(\Gamma(X',\psi^*(\mathcal N^\bullet)))\ar[d]\\
H^p(X',\psi^*(\mathcal F))
\otimes_{\Gamma(X',\psi^*(\mathcal O_X))}
H^q(X',\psi^*(\mathcal G))\ar[r]
&H^{p+q}(X',\psi^*(\mathcal H))
}
\]
Enfin, par fonctorialité, on a un diagramme commutatif
\[
\label{0.12.1.5.6-fr}
\tag{12.1.5.6}
\xymatrix@C=7pt{
H^p(X',\psi^*(\mathcal F))
\otimes_{\Gamma(X',\psi^*(\mathcal O_X))}
H^q(X',\psi^*(\mathcal G))\ar[r]\ar[d]
&H^{p+q}(X',\psi^*(\mathcal H))\ar[d]\\
H^p(X',f^*(\mathcal F))\otimes_{A'}H^q(X',f^*(\mathcal G))\ar[r]
&H^{p+q}(X',f^*(\mathcal H))
}
\]
et par combinaison des trois diagrammes
(\hyperref[0.12.1.5.4-fr]{12.1.5.4}),
(\hyperref[0.12.1.5.5-fr]{12.1.5.5}) et
(\hyperref[0.12.1.5.6-fr]{12.1.5.6}), on obtient le diagramme commutatif
(\hyperref[0.12.1.5.2-fr]{12.1.5.2}) cherché.
\end{env}

\begin{remark}[12.1.6]
\label{0.12.1.6-fr}
